\documentclass[11pt,a4paper]{article}
\usepackage[margin=25mm]{geometry}
\usepackage[T1]{fontenc}
\usepackage{lmodern,microtype,amsmath,amssymb,amsthm,booktabs,array,tabularx}
\usepackage[hidelinks]{hyperref}
\usepackage{url}
\newtheorem{theorem}{Theorem}
\newtheorem{lemma}[theorem]{Lemma}
\newtheorem{proposition}[theorem]{Proposition}
\theoremstyle{definition}
\newtheorem{definition}[theorem]{Definition}
\newcommand{\AC}{\{\mathrm A,\mathrm C\}}
\newcommand{\WC}{\mathcal W}
\newcommand{\score}{\operatorname{st}}
\newcommand{\comp}{\mathrel{\sim}}
\newcommand{\U}{\mathcal U_{14}}
\newcolumntype{Y}{>{\raggedright\arraybackslash}X}
\setlength{\emergencystretch}{1.5em}
\title{Twelve two-base-pair helices: an exact designability bound\protect\\
for a fixed family in the maxStacks model}
\author{Tatsuhiko Takimoto}
\date{25 October 2026 \quad Version 1.2}
\begin{document}
\maketitle
\begin{abstract}
Fix a family of noncrossing RNA targets with one enclosing two-base-pair
helix, $D-1$ inner two-base-pair helices, and one unpaired nucleotide in
each hairpin and between consecutive helices around the central loop.
In the maxStacks model with GU pairs allowed, no isolated pairs, and
minimum pair separation $j-i\geq2$, the largest designable value of $D$
is $12$. A 71-nucleotide sequence makes the $D=12$ target the unique
stack maximizer. Nonexistence for $D=13$ follows from a nucleotide-level
CNF and one DRAT refutation, also checked after conversion to LRAT
by a verified CakeML checker. For every $D\geq14$, a family-wise embedding
reduces a putative design to a 97-position relaxation whose 34 finite
cases are inconsistent. We give the normalization and local replacement
arguments connecting these certificates to the model, and separate
their verification obligations. The one-nucleotide hairpin convention
is essential to the stated target family; the result does not concern
the usual minimum of three unpaired hairpin nucleotides.
\end{abstract}

\section{The question}
Boury, Bulteau and Ponty~\cite{BBP} formulate inverse folding under stack
maximization and show how sufficiently long helices permit design at
large loop degree. Their Lemma~1 excludes a loop with degree $\Delta$
and equal helix length $h$ when $h<\log_6\Delta$. At $h=2$, this gives
the necessary condition $\Delta\leq36$; it does not assert that degree
$36$ is attainable. We are grateful for their model and degree--length
obstruction, which supply the question studied here.

We determine the maximum for one fixed family with helices of length
two and specified unpaired positions. Under the model below, a target
with twelve such helices admits a unique optimal fold, whereas every
target in this family with thirteen or more does not. The contribution
is this exact value for the fixed infinite family. The model, local
exchange obstructions and the use of finite certificates are not
claimed as new methods.

The proof has three parts: an explicit design at $D=12$, a direct
refutation at $D=13$, and a separate argument for all $D\geq14$.
The last part is needed because failure at $D=13$ alone does not imply
failure at larger $D$. Local replacements explain the finite reductions;
the certificate for $D=13$ can be checked without trusting the auxiliary
graph and context encodings used to construct it.

\section{Model, target family and theorem}\label{sec:model}
Let $\Sigma=\{\mathrm A,\mathrm C,\mathrm G,\mathrm U\}$ and let the
symmetric compatibility relation $\comp$ consist of
\[
 \mathrm{AU},\ \mathrm{UA},\ \mathrm{CG},\ \mathrm{GC},\
 \mathrm{GU},\ \mathrm{UG}.
\]
A structure $P$ on positions $0,\ldots,n-1$ is a set of pairs $(i,j)$,
$i<j$, with disjoint endpoints. For any two pairs $(i,j),(k,l)\in P$,
the crossing $i<k<j<l$ is forbidden.
Every pair must satisfy $j-i\geq2$ and have a neighboring nested pair:
at least one of $(i+1,j-1)$ and $(i-1,j+1)$ also belongs to $P$.
Call a structure satisfying these requirements a \emph{legal structure};
in particular, it has no isolated pairs. A structure is compatible
with a sequence $w$ if $w_i\comp w_j$ for every $(i,j)\in P$.
Its score is
\[
 \score(P)=\bigl|\{(i,j)\in P:(i+1,j-1)\in P\}\bigr|.
\]
A helix is a maximal run of consecutive nested pairs. A helix of
$\ell$ pairs contributes $\ell-1$ stacks, so a two-base-pair helix
contributes one stack. The pair-separation condition permits a
hairpin with just one unpaired nucleotide. It is the theoretical
convention of~\cite[Section~2]{BBP}, not the usual RNA convention
requiring at least three.

For $D\geq2$, put $n_D=6D-1$ and $s_k=6k-3$ for $1\leq k\leq D-1$.
Define
\begin{equation}\label{eq:target}
 T_D=\{(0,n_D-1),(1,n_D-2)\}
 \ \cup\ \bigcup_{k=1}^{D-1}\{(s_k,s_k+4),(s_k+1,s_k+3)\}.
\end{equation}
Equivalently, its dot-bracket string is
\[
 \texttt{((.}\quad
 \underbrace{\texttt{((.)).}\ \cdots\ \texttt{((.)).}}_{D-1\text{ copies}}
 \quad\texttt{))}.
\]
Whitespace in this display is only typographic. The enclosing helix
and the $D-1$ inner helices give score $D$. Thus $D$ counts all helices,
including the enclosing one, and is the degree of the central loop.
The unpaired positions between branches will be called separators.

\begin{definition}
A design of $T_D$ is a sequence compatible with $T_D$ for which every
other compatible legal structure $P\ne T_D$ satisfies $\score(P)<D$.
Let $D_2$ be the largest $D\geq2$ for which such a sequence exists.
\end{definition}

\begin{theorem}\label{thm:main}
In the model of Section~\ref{sec:model}, $T_{12}$ is designable and
$T_D$ is not designable for every $D\geq13$. Consequently $D_2=12$.
\end{theorem}

\section{A design attaining twelve}\label{sec:lower}
The following sequence has length 71:
\begin{center}\small\ttfamily
GAAAGCCUAGGACCACCCGGACGACGAACAGUACAAUGAGUAACAUGACACUAAUAAAUAAUCGCAGCAUC
\end{center}
Two dynamic programs, one using free-interval and paired-boundary
states and one enumerating maximal helices, return optimum $12$ and
exactly one optimal structure, namely $T_{12}$. The bundled checker
\path{checkers/core.py} implements the maximal-helix recurrence below;
\path{checkers/check_finite.py} compares it with direct enumeration.

The maximal-helix decomposition explains both optimization and counting.
For an interval $[i,j]$, let $M_0(i,j)$ be the best score over all legal
structures on that interval, and let $M_1(i,j)$ be the best score when
the pair $(i,j)$ is forbidden. Retain the number of optimal structures
alongside each score. A structure either
leaves $i$ unpaired or starts with a maximal helix
\[
 H(i,k,\ell)=\{(i+r,k-r):0\leq r<\ell\},\qquad \ell\geq2.
\]
Require compatibility of all these pairs and
$k-i-2(\ell-1)\geq2$. Its score contribution is
\begin{equation}\label{eq:dp}
 (\ell-1)+M_1(i+\ell,k-\ell)+M_0(k+1,j).
\end{equation}
For $M_1(i,j)$, exclude $k=j$; the unpaired option uses
$M_0(i+1,j)$. Empty intervals have score zero and one structure.
The inner state in~\eqref{eq:dp} forbids extension of the chosen
helix, so a structure has exactly one decomposition. Counts multiply
over the two intervals and add over equal-score alternatives. This
counts structures rather than different parses of the same structure.

A separately written implementation of this decomposition also returns
$(12,1,T_{12})$. It was compared with direct enumeration of noncrossing
matchings, followed by the legality tests, on 22,485 short sequences.
The finite comparison supports the implementation; the unique
decomposition is the mathematical reason for the recurrence. No
general claim about designability at all smaller $D$ is needed.

\section{Normalization and local obstructions}\label{sec:local}
Write $\WC=\{\mathrm{AU},\mathrm{UA},\mathrm{CG},\mathrm{GC}\}$ for
Watson--Crick pairs, and $\overline a$ for the unique Watson--Crick
complement of $a$.

\begin{lemma}[Normalization]\label{lem:normal}
If $T_D$ has a design, it has one whose target pairs all lie in $\WC$
and whose target-unpaired positions all lie in $\AC$.
\end{lemma}
\begin{proof}
The compatible-neighbor sets satisfy
$N(\mathrm A)=\{\mathrm U\}\subset N(\mathrm G)=\{\mathrm C,\mathrm U\}$
and
$N(\mathrm C)=\{\mathrm G\}\subset N(\mathrm U)=\{\mathrm A,\mathrm G\}$.
Replacing a single G by A, or U by C, cannot create a compatible pair.
Apply these replacements at target-unpaired positions. At each target
GU or UG pair, replace its G by A, making it a Watson--Crick pair.
Each position belongs to at most one target pair, so the changes can be
made sequentially while preserving the target. They only remove
compatible competitors. The target therefore remains uniquely optimal.
\end{proof}
This is an existence-preserving transformation, not a condition that
every original design must already satisfy. The global permutation
$\tau:\mathrm A\leftrightarrow\mathrm C$, $\mathrm G\leftrightarrow\mathrm U$
preserves compatibility, $\WC$ and $\AC$. We use it either to fix one
separator or to choose an orbit representative, never to impose both
restrictions without an additional argument.

Call a legal structure different from $T_D$ with score at least $D$
a \emph{rival}. An incompatible pair in every rival is necessary for a
design. Two replacement operations provide many such obstructions.

\begin{lemma}[Swap and three-pair replacement]\label{lem:replace}
Let two inner branches start at $s<t$.
Replacing their four pairs by
\begin{equation}\label{eq:swap}
 (s,t+4),\ (s+1,t+3),\ (s+3,t+1),\ (s+4,t)
\end{equation}
gives a rival. Replacing the enclosing helix and the branch at $s$ by
\begin{equation}\label{eq:outer-swap}
 (0,s+1),\ (1,s),\ (s+3,n_D-1),\ (s+4,n_D-2)
\end{equation}
also gives a rival.
Alternatively, choose consecutive three-position blocks starting at
$x<y$, each within the seven-position window of one of two branches
or at an enclosing-helix end. If $y-x\geq4$, replace the two original
helices by
\begin{equation}\label{eq:triplet}
 K(x,y)=\{(x,y+2),(x+1,y+1),(x+2,y)\}.
\end{equation}
For the choices specified below, this is a rival as well.
\end{lemma}
\begin{proof}
For~\eqref{eq:swap}, $t-s\geq6$, so the shortest pair has separation
$t-s-4\geq2$. The four pairs are nested and form two stacks.
Branches between $s$ and $t$ lie inside the innermost new pair;
the remaining branches are disjoint from the new pairs or enclose them.
In~\eqref{eq:outer-swap}, the two new helices are disjoint and enclose
the branches before and after $s$, respectively. Their shortest pairs
also have separation at least two.

In~\eqref{eq:triplet}, the innermost separation is $y-x-2\geq2$.
The new helix has three pairs and two stacks, replacing two stacks.
A branch window contains only that branch's paired positions and
target-unpaired positions. Hence none of the retained target pairs
uses a new endpoint. For the inner and enclosing choices below, all
retained branches lie inside a new pair or in a disjoint interval.
There is no crossing or isolated pair. Each replacement changes the
target and preserves at least its score; additional stacks can only
increase that score.
\end{proof}

\subsection{A graph on 32 words}
A normalized inner branch is the five-letter word
\[
 v=(a,b,h,\overline b,\overline a),\qquad a,b\in\Sigma,\quad h\in\AC.
\]
There are 32 words. Join distinct words when their $(a,b)$ types agree,
or when some internal consecutive triplet in one is reverse-compatible
with an internal consecutive triplet in the other. Reverse-compatible
means that all three opposing letter pairs are compatible.
This graph has 66 edges, obtained by evaluating these finite predicates.
Repeated types permit~\eqref{eq:swap}. For internal triplets, offsets
from branch start lie in $\{0,1,2\}$, so their start separation is at
least $6-2=4$, permitting~\eqref{eq:triplet}.
Thus the $D-1$ inner words of any normalized design form an independent
set, with no repeated word or type.

\subsection{The neighboring separators}
Extend a word to its seven-letter context
\[
 L\ a\ b\ h\ \overline b\ \overline a\ R,\qquad L,R\in\AC.
\]
The consecutive triplets start at offsets $-1,0,1,2,3$ from the
position of $a$. Two contexts with offsets $r,q$ satisfying
$|r-q|\leq2$ admit a legal three-pair replacement in either branch
order. For a single window, the six legal fragments of score at
least one other than the original \texttt{.((.)).} are
\[
\begin{gathered}
 \texttt{((.))..},\quad\texttt{..((.))},\quad\texttt{((..)).},\\
 \texttt{.((..))},\quad\texttt{((...))},\quad\texttt{(((.)))}.
\end{gathered}
\]
Each replaces that branch while the other target pairs remain.

For the enclosing helix, the end triplets occupy $[0,2]$ and
$[n_D-3,n_D-1]$. Its context consists of the four paired letters and
the two separators at $2$ and $n_D-3$. Besides the outer
swap~\eqref{eq:outer-swap}, its relevant replacements are
\[
\begin{array}{ll}
 \text{left triplet:} & K(0,s+r),\quad r\in\{1,2,3\},\\
 \text{right triplet:} & K(s+r,n_D-3),\quad r\in\{-1,0,1\},\\
 \text{first shift:} & \{(0,n_D-2),(1,n_D-3)\},\\
 \text{second shift:} & \{(1,n_D-1),(2,n_D-2)\}.
\end{array}
\]
Here $s=6j-3$ is an actual inner-branch start. The triplet moves remove
the enclosing helix and that branch; a shift removes just the enclosing
helix. In each case all other target pairs are retained.

The actual contexts share separators between consecutive branches.
Allowing their copies to vary independently relaxes the necessary
conditions. A solution of this relaxation need not be a design.
Its inconsistency, together with the embedding below, excludes a design.

\section{Excluding every \texorpdfstring{$D\geq14$}{D >= 14}}\label{sec:upper}
\subsection{A relaxation on 97 positions}
Let the auxiliary target have dot-bracket string
\[
 \texttt{((}\ \texttt{.}\
 \underbrace{\texttt{.((.)).}\ \cdots\ \texttt{.((.)).}}_{13\text{ copies}}
 \ \texttt{.}\ \texttt{))}.
\]
It has 97 positions, enclosing pairs $(0,96),(1,95)$, external
separators at $2,94$, and thirteen seven-position blocks
$3+7k,\ldots,9+7k$ for $0\leq k\leq12$.
The auxiliary branch starts are $u_k=4+7k$; a block reads
$L\,a\,b\,h\,\overline b\,\overline a\,R$.
All target pairs are Watson--Crick and all target-unpaired letters are
in $\AC$, with no separator fixed by color symmetry.

Define $\U$ to be the CNF requiring these conditions and forbidding the
1,731 rivals in Table~\ref{tab:families}. All unspecified target pairs
are retained. Nucleotide variables are one-hot, and a variable $p_{ij}$
is equivalent to compatibility of the two letters, using the encoding
in Section~\ref{sec:cnf}. Here $p_{ij}$ is introduced for every $i<j$,
including adjacent positions. This is just a compatibility variable;
every rival still obeys $j-i\geq2$.
The formula has 5,044 variables and $67,129+1,731=68,860$ clauses.

\begin{table}[ht]
\centering\small
\begin{tabularx}{\linewidth}{@{}>{\raggedright\arraybackslash}p{26mm}Yr@{}}
\toprule
Rival family & Replacement & Count\\
\midrule
One window & Replace its branch by any of the six other legal
seven-position fragments with score at least one. & 78\\
Inner swap & Use~\eqref{eq:swap} with $(s,t)=(u,v)$ for two auxiliary branches. & 78\\
Outer swap & Use~\eqref{eq:outer-swap} with $s=u$ and $n_D=97$. & 13\\
Inner triplets & $K(u+r,v+q)$, $u<v$, $r,q\in\{-1,0,1,2,3\}$,
$|r-q|\leq2$. & 1,482\\
Left outer triplet & $K(0,u+r)$, $r\in\{1,2,3\}$. & 39\\
Right outer triplet & $K(u+r,94)$, $r\in\{-1,0,1\}$. & 39\\
Outer shift & Replace the enclosing pairs by
$\{(0,95),(1,94)\}$ or $\{(1,96),(2,95)\}$. & 2\\
\bottomrule
\end{tabularx}
\caption{The seven families of auxiliary rivals, with auxiliary branch
starts $u,v\in\{4+7k:0\leq k\leq12\}$. A triplet replacement
removes the two original helices it uses. Each structure is legal,
differs from the auxiliary target, and has score at least 14.}
\label{tab:families}
\end{table}

\subsection{Covering the finite possibilities}\label{sec:cover}
The thirteen words have distinct $(a,b)$ types and form an independent
set in the 32-vertex graph. There are
\[
 \binom{16}{13}2^{13}=4,587,520
\]
type-and-hairpin choices before the graph restrictions. An exhaustive
enumeration leaves 68 sets of thirteen words. Each rejected choice
contains a forbidden pair supported by an actual swap or internal
triplet rival. The retained sets split into 34 two-element orbits under
$\tau$.

To pass from word sets to fixed cases, one must also account for their
order. Each of the twelve adjacent permutations of the seven-position
blocks, and the letter permutation $\tau$, preserves the multiset of
\emph{all clauses} of $\U$. This was checked by literal substitution.
The adjacent permutations generate every block permutation. Therefore
any model can be transformed to one with words in lexicographic order
(using $\mathrm A<\mathrm C<\mathrm G<\mathrm U$)
and belonging to one of the 34 representatives. Their thirteen-word
sets are listed in \path{data/U14-cases.json}. This is a symmetry of
the selected formula, not an assertion about arbitrary RNA structures.

Each case adds 65 nucleotide units fixing the five letters in each of
thirteen branches. Its CNF has 68,925 clauses. All 34 cases are
inconsistent. In addition to certificate acceptance by the public
DRAT checker, a separate unit-propagation implementation produces a
contradiction in every case, and a replay checks the reason clause for
each propagation. Thus this finite step can be checked using the case
inputs alone, without the supplied DRAT additions.

\begin{proposition}\label{prop:u}
The formula $\U$ is unsatisfiable.
\end{proposition}
\begin{proof}
Every model would have one of the 68 surviving word sets. Formula
symmetry carries it to one of the 34 fixed cases, contradicting that
case's checked inconsistency.
\end{proof}

\subsection{Returning a rival to the original target}\label{sec:embedding}
\begin{lemma}\label{lem:embedding}
A design of $T_D$ for any $D\geq14$ would give a model of $\U$.
\end{lemma}
\begin{proof}
Normalize the design and choose any thirteen inner branches, indexed
$j_0<\cdots<j_{12}$. Copy the letters at actual positions
$0,1,n_D-2,n_D-1$ to auxiliary positions $0,1,95,96$ and those at
$2,n_D-3$ to $2,94$. For each block set
\begin{equation}\label{eq:map}
 \phi(3+7k+r)=6j_k-4+r,\qquad 0\leq r\leq6.
\end{equation}
Assign each auxiliary letter the letter at its image. The target and
unpaired-letter clauses hold, and assign every $p_{ij}$ its actual
compatibility truth value.

The map $\phi$ is not injective on all positions: consecutive selected
branches can share a separator, and the first or last block can share
a separator with an enclosing end. We check that every listed rival
still gives a valid replacement in the original structure.

For a one-window rival, the seven positions stay consecutive and
distinct. Its endpoints outside the branch were target-unpaired, so
all other target helices may be retained. For an inner triplet rival,
actual branch starts $s=6j_a-3<t=6j_b-3$ satisfy $t-s\geq6$. Hence
\[
 (t+q)-(s+r)\geq6-2=4.
\]
The two triplets have distinct endpoints and legal separation. The
potential shared-separator collision would require the left offset
$r=3$ and right offset $q=-1$, which is excluded by $|r-q|\leq2$.
Retained branches between the chosen ones are enclosed by the new
helix; those outside do not cross it.

For a left outer triplet, $s\geq3$ and $r\geq1$ give
$s+r\geq4$. For a right outer triplet, $s\leq n_D-8$ and $r\leq1$
give $(n_D-3)-(s+r)\geq4$. These same restrictions prevent use of a
shared separator twice. Retained branches lie inside the new helix or
in a disjoint interval.

Inner swaps use only paired positions, with shortest new separation
$t-s-4\geq2$. Outer swaps have shortest pairs $(1,s)$ and
$(s+4,n_D-2)$, again of separation at least two. Their nesting and
disjointness follow the proof of Lemma~\ref{lem:replace}.
For an outer shift, the enclosing endpoints and external separator
remain consecutive actual positions, so the shifted helix remains
legal and encloses every retained inner branch.

In every family, the replaced one or two target stacks are replaced
by at least as many stacks. Add all untouched target helices, including
unselected branches. The resulting structure is noncrossing, has no
isolated pair, differs from $T_D$, and has score at least $D$.
If the auxiliary rival were compatible, the same letter pairs would
make this actual rival compatible, contradicting the design. Thus all
rival-exclusion clauses of $\U$ hold, giving the asserted model.
\end{proof}
Proposition~\ref{prop:u} and Lemma~\ref{lem:embedding} exclude all
$D\geq14$. Tests of finitely many embeddings supplement this argument;
the inequalities and family-wise construction establish its universal
quantifier. This route uses an auxiliary target, not a direct single
refutation for the canonical $T_{14}$.

\section{Excluding \texorpdfstring{$D=13$}{D = 13} by one pure CNF}\label{sec:d13}
\subsection{Nucleotides, compatibility and rivals}\label{sec:cnf}
For the 77 positions of $T_{13}$, introduce one-hot variables
$x_{i,a}$ for $a\in\Sigma$ and variables $p_{ij}$ for $j-i\geq2$.
Require Watson--Crick target pairs and A/C at target-unpaired positions.
Use $\tau$ once to fix the first separator, position 2, to A.

For each compatible ordered letter pair $a\comp b$, add
$\neg x_{i,a}\lor\neg x_{j,b}\lor p_{ij}$. For each $a\in\Sigma$,
also add
\begin{equation}\label{eq:support}
 \neg p_{ij}\lor\neg x_{i,a}\lor\bigvee_{b\in N(a)}x_{j,b},
\end{equation}
and its endpoint-reversed version. The six forward clauses and eight
support clauses, with one-hot letters, make $p_{ij}$ equivalent to
actual compatibility. Redundant units set target $p_{ij}$ true and
target-unpaired/target-unpaired $p_{ij}$ false.

Let $\mathcal R_{13}$ be the supplied list of 227,280 distinct rivals.
For every $P\in\mathcal R_{13}$, add
\begin{equation}\label{eq:rival-clause}
 \bigvee_{(i,j)\in P\setminus T_{13}}\neg p_{ij}.
\end{equation}
Every listed structure has been checked for balanced brackets,
disjoint endpoints, noncrossing, minimum separation, absence of isolated
pairs, difference from the target, and score at least 13.
There are 227,255 rivals of score 13 and 25 of score 14.
The list need not contain every legal structure: any design must avoid
every member of any such list.

The resulting formula $F$ has 3,158 variables: 308 nucleotide variables
and 2,850 compatibility variables. Its 41,128 base clauses comprise
539 one-hot clauses, 312 target Watson--Crick exclusions, 50 unpaired
restrictions, one separator unit, 39,900 compatibility clauses, and
326 redundant units. Adding the rivals gives 268,408 clauses.
The input has no word, set or context variables.

\begin{proposition}\label{prop:d13}
The formula $F$ is unsatisfiable, and $T_{13}$ has no design.
\end{proposition}
\begin{proof}
Unsatisfiability is certified by the single checked DRAT stream
identified in Appendix~\ref{app:cert}. A hypothetical design could be
normalized by Lemma~\ref{lem:normal}, then color-swapped to fix position
2. Assign its letters and their compatibility values to $x$ and $p$.
All base clauses would hold. Since it is a design, each legal rival
has an incompatible pair, which cannot be a target pair; hence
\eqref{eq:rival-clause} holds. This would satisfy $F$, a contradiction.
\end{proof}

\subsection{How the refutation is assembled}
The graph and context reductions organize the certificate's construction.
There are 1,267 independent word sets of size twelve. Context
obstructions dispose of 1,123, and position-dependent cases dispose of
the remaining 144. In the final certificate, word and context indicators
are introduced as fresh definitions within the proof. Case conclusions
are guarded by the negation of their conditions, so no case condition
becomes an assumption of the input $F$.

Prefix-occurrence variables record which words occur in the first $k$
branch positions. The proof rules out the final word sets and propagates
these exclusions backwards to the empty prefix, ending in the empty
clause. Definitions are checked by the resolution asymmetric tautology
(RAT) rule; subsequent reverse unit propagation (RUP) steps and
deletions are part of the same stream. Its maximum variable index is
8,727, although $F$ itself has only 3,158 variables. The checker verifies
the additions against the original input and preceding proof state.
Consequently, neither the older auxiliary encodings nor an assertion
that the 1,267-set list is complete is an extra premise of
Proposition~\ref{prop:d13}.

The design in Section~\ref{sec:lower}, Proposition~\ref{prop:d13}, and
the universal exclusion in Section~\ref{sec:upper} prove
Theorem~\ref{thm:main}.

\section{Computer assistance and proof obligations}\label{sec:obligations}
Table~\ref{tab:obligations} separates the arguments that must be read
mathematically from the finite objects checked by programs.
For $D=13$, the verification route is to rebuild $F$, check its rivals,
and check the refutation. For $D\geq14$, it is to check the 34-case
cover and inconsistencies, then the general embedding. The local
lemma chain explains why these finite objects arise.

\begin{table}[ht]
\centering\small
\begin{tabularx}{\linewidth}{@{}>{\raggedright\arraybackslash}p{28mm}YY@{}}
\toprule
Obligation & Prose argument & Finite or certificate check\\
\midrule
Attainment at 12 & Unique maximal-helix decomposition & Two DP
recurrences give optimum 12 and count 1; separate implementation agrees.\\
Normalization & Neighbor-set inclusion preserves the target and removes
competitors & Alphabet relations and color symmetry checked on all letters.\\
Local replacements & Endpoint geometry, noncrossing and stack lower
bound & Word predicates; all listed rivals tested from the definition.\\
All $D\geq14$ & Family-wise embedding, including shared separators &
Rebuild auxiliary CNF; 68 word sets, 34 orbits, full-clause symmetries;
34 contradictions, with independent unit-propagation replay and
additional LRAT checks by \texttt{cake\_lpr}.\\
$D=13$ & A normalized design would satisfy $F$ & Rebuild all 268,408
clauses, validate 227,280 rivals, check the DRAT refutation of $F$
and its LRAT conversion with \texttt{cake\_lpr}.\\
\bottomrule
\end{tabularx}
\caption{Responsibility for the computer-assisted theorem. A checker
acceptance alone does not supply the model-to-formula argument.}
\label{tab:obligations}
\end{table}

The accepted refutation is stored as seven ordered gzip files. Their
concatenated decompressed stream is binary DRAT, with
2,049,035,259 bytes; the compressed total is 613,931,594 bytes.
Verification uses the public \texttt{drat-trim} checker~\cite{DRAT},
with \texttt{-i -f}: explicit binary input from standard input and
forward verification. The independent run builds with assertions
enabled and returns \texttt{VERIFIED}, exit code zero, on the original
$F$. Appendix~\ref{app:cert} fixes the input, rivals and raw proof by
SHA-256. The companion inventory specifies every part and its order.

\subsection{Checking the LRAT with a verified program}
The original $F$ and all 34 auxiliary case CNFs also yield the exact
success line \texttt{s VERIFIED UNSAT} from the official, unmodified
ARMv8 build of \texttt{cake\_lpr}~\cite{CAKE}, using LRAT proofs converted
from the DRAT certificates. The $D=13$ LRAT has 5,428,472,887 raw bytes
and is stored in a 1,598,317,453-byte gzip file. A separate build replays
the stored certificates without rerunning their converter and obtains
the same success result. Corrupted-proof controls are rejected.

The HOL4 theorem \texttt{machine\_code\_sound} for the pinned CakeML
revision connects the compiled checker to its parsed CNF semantics.
In the usual two-argument CNF/proof mode, under its ARM execution,
foreign-function-interface and filesystem assumptions, the success
output implies that the parsed CNF is unsatisfiable. Successful process
exit alone is insufficient. The DRAT-to-LRAT converter, including its
local streaming changes, need not be trusted for this implication:
\texttt{cake\_lpr} checks the resulting proof against the input CNF.

The remaining trust includes HOL4's logic and kernel, the formal
semantics, the correspondence between the ARM model and the machine,
the C input/output interface, assembler/linker and operating system,
and the provenance linking the distributed assembly to the proved
compilation result. We used the distributed ARM assembly; we did not
rebuild HOL4/CakeML or regenerate it from the proof. The normalization,
rival legality, CNF/model correspondence, 34-case cover, general
embedding and design DP remain the separate obligations in
Table~\ref{tab:obligations}. The checker theorem does not itself
formalize those RNA arguments.

\section{Independent verification}\label{sec:independent}
A separate verification implementation used the statement, model and
CNF specifications, the auxiliary embedding description, the rival
lists and the certificate manifest. It did not use the producer's
generation code, auxiliary proof-lifting explanation, prior check
results or the coordinator's implementation. Public checker
documentation and parser source were available to it. Audit details,
dates and responsibility are recorded once in the companion
\texttt{VERIFICATION-RECORD.md}.

This implementation recovered the 71-base design result, checked every
rival, and regenerated both input CNFs byte for byte. It regenerated
all seven auxiliary rival families, checked the entire word-set cover
and the full-formula symmetries, and reproduced all 34 case
contradictions. It checked the seven-part proof's hashes and complete
raw stream, and accepted the $D=13$ refutation. It also reviewed the
embedding family by family. Tests of 50,199 mapped rivals supplement
that general argument.

The three initial DRAT checks use the same public \texttt{drat-trim}
implementation. The additional \texttt{cake\_lpr} checks and replay
supply a different checker with a proved correctness theorem, but are
not part of the separate model-verification implementation described
above: that implementation used \texttt{drat-trim} and did not rerun
\texttt{cake\_lpr}. The verification record distinguishes the roles.
For the auxiliary cases, \texttt{drat-trim} accepts by unit propagation
on the inputs, so this does not mean every supplied DRAT addition was
exercised. Separate propagation traces and replay check that finite
step as well. Normalization, the dynamic-programming decomposition and
the embedding remain prose arguments. The full RNA theorem has not
been formalized in a proof assistant, and this article has not undergone
human peer review.

\section{Scope}
The maximum $12$ concerns precisely~\eqref{eq:target}, with every
separator and hairpin length fixed as stated. It gives no maximum
degree for arbitrary RNA trees, nor a classification of every
$D\leq12$. Changing the separator length can change the bound: the
variant of~\eqref{eq:target} with all separators removed and one-base
hairpins retained admits a design at $D=25$.\footnote{\raggedright The 124-nucleotide sequence
{\ttfamily\footnotesize
ACUCAGAUCAGACGCCGAGACUAGACUCCAG\allowbreak{}
GCCAGGGGACCGGACCUUAAAACCGUACCGU\allowbreak{}
UGCCAUGCCAGACUCGACUCAUCAUAUCAUU\allowbreak{}
ACUAUACUACUCAGCUCAGAAAUUGCCGCGU}
designs the target \texttt{((} followed by 24 copies of \texttt{((.))} and then \texttt{))}. This example can be checked with any maxStacks dynamic program implementing the model of Section~\ref{sec:model}. It is outside this article's claims and certificates.} This variant is being
studied separately and is outside the claims and certificates of the
present paper. If hairpins must contain three unpaired nucleotides, these
targets themselves cease to be legal; lengthening them defines a
different family. No conclusion about the Turner energy model,
physical molecules or experimental folding follows from this theorem.
The exact value is the proposed contribution within the reviewed
literature; a priority claim is not established by a negative search.

\paragraph{Use of AI.}
Codex and Claude assisted the mathematical arguments, search, software,
verification implementations, literature comparison and drafting.
Tatsuhiko Takimoto is responsible for the manuscript. The computational
and mathematical assurances are those stated above.

\begin{thebibliography}{9}
\raggedright
\bibitem{BBP}
Th\'eo Boury, Laurent Bulteau and Yann Ponty.
RNA Inverse Folding Under Stacked Base Pair Maximization.
In \emph{26th International Conference on Algorithms for Bioinformatics
(WABI 2026)}, LIPIcs 390, 12:1--12:21, 2026.
\href{https://doi.org/10.4230/LIPIcs.WABI.2026.12}{doi:10.4230/LIPIcs.WABI.2026.12}.
\bibitem{DRAT}
Marijn J. H. Heule and contributors.
\emph{The DRAT-trim proof checker}. Public software and documentation.
\url{https://github.com/marijnheule/drat-trim}.
The exact revision and build used for verification are recorded in the
companion verification record.
\bibitem{CAKE}
cake\_lpr contributors and the CakeML project.
\emph{LPR proof checking with CakeML}.
\url{https://github.com/tanyongkiam/cake_lpr};
\emph{Verified proof checking}, \url{https://cakeml.org/checkers.html}.
ARM correctness theorem:
\href{https://github.com/CakeML/cakeml/blob/fb377b4bb704497c921cde68ccc8da3b4f0e9132/examples/lpr_checker/array/compilation/proofsARM8/lpr_arrayProofARM8Script.sml}{\texttt{lpr\_arrayProofARM8.machine\_code\_sound}}.
\end{thebibliography}

\appendix
\section{Certificate identifiers}\label{app:cert}
The following hashes identify the mathematical inputs and accepted
proof, independently of local filenames. A rival-list hash refers to
the uncompressed ASCII dot-bracket lines, including their newlines.
The accompanying package specifies variable numbering, clause order and
literal order, and rival-list order in \path{claims/SPEC.md} and
\path{claims/RIVAL-UPPER.md}. The thirteen-word sets for all 34 cases,
and the CNF and LRAT hashes, are in \path{data/U14-cases.json}.
The names, order and individual hashes of the seven compressed $D=13$
DRAT parts, and the large-proof inventory, are in
\path{data/large-proofs.json}. The large proofs are distributed separately
as release attachments. The article text alone is not sufficient for
byte-exact regeneration; the serialization details are supplied in those
package files.

\begingroup\small
\noindent\textbf{$D=13$ input $F$:} 3,158 variables; 268,408 clauses;
20,535,942 bytes.\\
\path{4f8100bbe4362bf45db9b4be068eb115bf1c48c3f556ce5c3bceff1e7f5a4a6c}

\medskip\noindent\textbf{$D=13$ rival list:} 227,280 structures.\\
\path{e1a149ce72fc055faa1e1f459ed1fd2f885df3d1c2795dab490154ac687f042a}

\medskip\noindent\textbf{$D=13$ complete raw binary DRAT:}
2,049,035,259 bytes.\\
\path{ef0a5a3cd19984a3f8973740d247ea229c57044a39d8b4068895e3493b0aa116}

\medskip\noindent\textbf{$\U$ input:} 5,044 variables; 68,860 clauses.\\
\path{5ff81ea47c5fa1f66300587ea0d036cae32598c2070374e8ee423056f89f715f}

\medskip\noindent\textbf{Auxiliary rival list:} 1,731 structures.\\
\path{0ad72ae92ae2ab1f141d6a4ed213a0ac6df72cd154430fdccae30a64e370374c}

\medskip\noindent\textbf{$D=13$ binary LRAT, gzip:}
1,598,317,453 bytes.\\
\path{8705abc77e86aa52856ef6c36f4bc01037693e7ef57fda96e0bf48f56cbc523a}

\medskip\noindent\textbf{Expanded binary LRAT:} 5,428,472,887 bytes.\\
\path{9b54d5df00f8db04a2d9d51688175ed75df683e54839c28db56fda714fddbd91}

\medskip\noindent\textbf{Official \texttt{cake\_lpr} revision:}\\
\path{a36874a8b750b43fe4b385b8ddbf5b033e46a3fa}

\medskip\noindent\textbf{ARMv8 checker executable SHA-256:}\\
\path{2fee767af4478ed817f8bdf6feea5866e5e59f1f67c2a4477e98ee961612881d}
The executable hash identifies this build, not every build on other
systems. The companion record pins HOL4 and CakeML revisions and
documents assembly and C-wrapper provenance.
\endgroup

\medskip
For a local check, concatenate the decompressed proof parts in the
inventory order and feed the resulting binary stream to
\texttt{drat-trim F.cnf -i -f}. An input hash, a proof hash and a checker
acceptance establish a propositional refutation of that input; the
normalization, rival legality and model correspondence must also be
checked as specified in Table~\ref{tab:obligations}.

To check the stored LRAT, decompress its gzip to the standard input of
\texttt{cake\_lpr}, with \texttt{F.cnf /dev/stdin} as the two checker
arguments. The C wrapper used here accepts heap and stack options
\texttt{--CML\_HEAP\_SIZE=4096} and \texttt{--CML\_STACK\_SIZE=128} before
those arguments. Require the exact \texttt{s VERIFIED UNSAT} line and
normal termination. Each auxiliary case must use its own CNF and
corresponding LRAT, as identified in the inventory.
\end{document}
